% =====================================================================
%  CHAPTER 8: 儿童少年常见病（对应大纲第七章）
%  参考PPT：第八章 儿童少年常见疾病（邹志勇，北京大学）
% =====================================================================
\chapter{儿童少年常见病}
\setcounter{chapter}{8}
\chapterintro{
掌握儿童少年常见病的概念和类型；掌握近视的定义、分类、流行特点、病因（遗传与环境因素）及防控措施（难点*）；掌握单纯性肥胖的定义、流行特点、病因（遗传、个体、家庭、社会环境因素）及综合防控（难点*）；掌握营养不足与缺铁性贫血的流行特点、发生原因及预防控制；掌握脊柱弯曲异常的流行特点、发生原因及预防控制（难点*）；掌握龋齿和牙周病的流行特点、四联致病因素及预防控制。
}

\begin{defbox}
\textbf{儿童少年常见病（Common Diseases in Children and Adolescents）/ 学生常见病：}是指在儿童少年中发生率较高的一类疾病，其发病与其正处于生长发育阶段有关，亦受到幼儿园和学校集体生活环境的影响。

\textbf{常见病类型：}近视、肥胖、营养不足与营养缺乏病、脊柱弯曲异常、心理健康问题、龋齿和牙周病、哮喘等。
\end{defbox}

% ----- 8.1 近视 -----
\section{儿童少年近视（难点*）}

\subsection{近视流行与危害}

\begin{keybox}
\textbf{（一）定义和分类}

\textbf{近视：}人眼在调节放松状态下，平行光线经眼球屈光系统后聚焦在视网膜之前，是屈光不正的一种类型。

\textbf{视力不良/视力低下：}各种原因导致的视力低于一定水平的总称，包括近视、远视和散光等屈光不正和弱视等其他眼病。

\textbf{等效球镜（SE）：}SE = 球镜度 + 1/2柱镜度。在流行病学调查中，使用最多的近视判定标准是SE $\leq$ $-$0.50D（即近视50度以上）。高度近视是SE $\leq$ $-$6.0D（即近视600度及以上）。

\textbf{分类：}
\begin{enumerate}[leftmargin=*]
    \item 按屈光成分：屈光性近视、轴性近视
    \item 按病程进展和病理变化：单纯性近视、病理性近视
    \item 按近视度数：低度近视、中度近视、高度近视（SE $\leq$ $-$6.0D）
    \item 按公共卫生层面防控策略分期：近视前驱期（近视前驱状态）、近视发展期、高度近视期、病理性近视期
\end{enumerate}

\textbf{（二）近视流行特点}

中国儿童少年近视患病率是世界上最高的国家之一，近视检出率随年龄、学龄而上升，学龄影响比年龄更大；城市高于乡村，但乡村学生的近视率增长更快；女童高于男童。

\textbf{（三）近视危害}

近视的典型症状是远视力下降。高度近视者发生视网膜脱离、撕裂、裂孔、黄斑出血、新生血管和开角型青光眼的危险性增高，严重者导致失明。
\end{keybox}

\subsection{近视病因与危险因素}

\begin{keybox}
\textbf{（一）遗传因素：}双生子和家系研究结果显示，近视遗传度为70\%$\sim$80\%，这可能与家系和同卵双生子中具有相似的环境有关。

\textbf{（二）环境因素：}\imp{户外活动时间是近视发生发展的重要保护性因素}，户外活动对近视的保护作用最重要的是户外时间长短，而非参加的具体运动类型，其因果关联得到随机对照试验研究的证实。

\textbf{（三）体质健康因素：}体质弱、健康状况差、早产儿、低出生体重儿等个体眼球壁发育不良，韧性不足，易致眼轴变长；\imp{青春期突增阶段是近视程度加重的关键期}，性发育越早的青少年，近视发展越快。
\end{keybox}

\subsection{近视预防控制}

\begin{keybox}
\begin{enumerate}[leftmargin=*]
    \item \imp{筛查视力不良与近视}：建立中小学生视力屈光筛查转诊技术流程。
    \item \imp{建立视力健康档案}：对0$\sim$6岁儿童和中小学生进行定期视力检查，建立视力健康档案，动态观察屈光状态发展变化，早期发现近视的倾向或趋势。
    \item \imp{培养健康用眼行为}：个体、家庭和学校应当积极培养"每个人都是自身健康第一责任人"的意识。
    \item \imp{建设视觉友好环境}：家庭、学校、医疗卫生机构、政府相关部门、媒体和其他社会团体等各界力量要主动参与。
    \item \imp{增加日间户外活动}：户外较强的光照强度可以促进视网膜多巴胺的释放，抑制眼轴的延长，延缓近视进展。户外具有开阔的视野，可以减轻眼调节负荷，放松睫状肌。
    \item \imp{规范视力健康监测与评估}：及时了解学生群体中视力不良、近视分布特点及变化趋势，确定高危人群及高危因素。
    \item \imp{分级诊疗与矫治}：有高危因素或远视储备不足者建议改变高危行为；裸眼视力下降者建议到医疗机构接受医学验光等屈光检查；已经近视者佩戴框架眼镜是矫正屈光不正的首选方法；使用低浓度阿托品或佩戴角膜塑形镜（OK镜）减缓近视进展时建议到正规医疗机构在医生指导下进行。
\end{enumerate}
\end{keybox}

% ----- 8.2 肥胖 -----
\section{儿童少年肥胖}

\subsection{肥胖流行与危害}

\begin{keybox}
\textbf{（一）定义和分类}

\textbf{肥胖：}是在遗传、环境因素交互作用下，因能量摄入超过能量消耗，导致体内脂肪积聚过多，从而危害健康的慢性代谢性疾病。

按病因可分为原发性肥胖和继发性肥胖；按全身脂肪组织分布部位分为腹型肥胖和周围型肥胖。\imp{原发性肥胖（单纯性肥胖）}的发生与遗传、饮食、身体活动水平、生活方式等有关，儿童肥胖大多属于此类。

\textbf{（二）肥胖流行特点}

全球：1975年全球儿童青少年（5$\sim$19岁）肥胖率不到1\%；到2022年时，有8\%儿童青少年存在肥胖，女性肥胖率增加到6.9\%（约6500万名），男性肥胖率增加到9.3\%（约9420万名）。

中国：7$\sim$18岁儿童青少年肥胖检出率由1985年的0.1\%增长至2019年的9.6\%，增长了75.6倍。乡村儿童青少年肥胖增长更为明显，预计2025年乡村儿童青少年超重肥胖检出率将全面超过城市，出现"城乡逆转"。

\textbf{（三）肥胖危害}

健康危害：空腹血糖升高、糖耐量受损、2型糖尿病、代谢综合征和脂肪肝等。心理-行为危害：可导致心理压抑，严重时造成心理精神障碍（如抑郁和自杀意念等）。社会危害：沉重的经济负担（肥胖控制治疗成本、生产力损失等）。
\end{keybox}

\subsection{病因与危险因素}

\begin{defbox}
\textbf{单纯性肥胖病因复杂}，存在遗传、个体和环境因素的影响及其交互和相互作用。

\begin{enumerate}[leftmargin=*]
    \item \imp{遗传因素}：影响肥胖发生、发展的重要因素，但不是决定因素，肥胖是多基因的关联、累加结果，属于多基因遗传。
    \item \imp{个体因素}：在肥胖遗传易感性存在情况下，95\%以上的肥胖发生与能量过度摄取和/或能量消耗不足相关的生活方式有关。因素包括膳食摄入、饮食行为、以及体力活动。
    \item \imp{家庭因素}：生命早期的环境可能会影响器官结构和功能，并影响日后的健康。
    \item \imp{社会环境因素}：传统的健康观念（如"胖是健康，以胖为福"）影响认知；全球化、城市化带来的膳食快餐化，媒体对饮食行为的影响等，共同构成肥胖易感环境。
\end{enumerate}

\textbf{肥胖筛查标准（中国标准）：}WGOC（中国肥胖问题工作组）标准——针对7$\sim$18岁儿童少年制定了分性别、年龄超重肥胖BMI筛查界值，18岁男女超重、肥胖筛查界值点分别确定为24kg/m$^2$和28kg/m$^2$，与成人标准对接。实际工作应用时，可以先按BMI标准筛查肥胖儿童少年，再用高腰围界值点区分腹型肥胖和周围型肥胖。
\end{defbox}

\subsection{肥胖综合防控}

\begin{enumerate}[leftmargin=*]
    \item \imp{强化家庭责任}：帮助儿童养成科学饮食行为；培养儿童积极身体活动习惯；做好儿童少年体重及生长发育监测；加强社区支持。
    \item \imp{强化学校责任}：办好营养与健康课堂；改善学校食物供给；保证在校身体活动时间。
    \item \imp{强化医疗卫生机构责任}：加强孕期体重管理；加强儿童少年体重管理；加强肥胖儿童少年干预。
    \item \imp{强化政府责任}：加强肥胖防控知识技能普及；强化食物营销管理；完善儿童少年体育设施。
\end{enumerate}

% ----- 8.3 营养不足和缺铁性贫血 -----
\section{儿童少年营养不足和缺铁性贫血}

\subsection{营养不足}

\begin{keybox}
\textbf{营养不足的表现：}
\begin{enumerate}[leftmargin=*]
    \item \imp{生长迟缓}：表现为年龄别身高低、发育迟缓会阻碍儿童发挥其身体和认知潜力
    \item \imp{消瘦}：表现为身高别体重低、中度或重度消瘦的年幼儿童面临更高死亡风险
    \item \imp{体重不足}：表现为年龄别体重低，体重不足的儿童可能会生长迟缓、消瘦或两者兼有
\end{enumerate}

\textbf{发生原因：}营养素摄入不足；疾病（肠道寄生虫感染、肺结核、肝炎等慢性疾病）；膳食结构和饮食习惯不良；早期生长潜力未充分发挥（6个月至3岁是身高增长最大潜能时期）；体像认知和心理因素（盲目节食追求"体型美"）。

\textbf{预防控制：}2017年国务院发布国民营养计划（2017—2030年），旨在降低5岁以下儿童的贫血率并控制肥胖率的上升趋势。针对中国儿童少年群体，应针对营养摄入相对不足、膳食结构不合理等，采取家庭、社区和学校相结合的营养知识、行为习惯、生活方式宣教。
\end{keybox}

\subsection{缺铁性贫血}

\begin{defbox}
\textbf{流行特点：}全球范围内，40\%的6至59月龄儿童受到贫血的影响。在中国，中国居民贫血患病率持续下降，但2019年中小学生贫血率反弹至11.1\%。

\textbf{发生原因：}
\begin{enumerate}[leftmargin=*]
    \item \imp{体内需求量大}：儿童少年生长发育旺盛，尤其青春期生长突增以及组织代谢的氧需求和女童月经周期等铁的需求量很大
    \item \imp{铁的额外丢失}：月经不调性失血（首位原因）、外伤、疾病（胃十二指肠溃疡、肠道蠕虫感染、慢性腹泻等）、环境污染（铅、苯、镉等中毒破坏红细胞）
\end{enumerate}

\textbf{预防控制：}去除病因（纠正不良饮食习惯和食物组成）；一般治疗（合理膳食、增加富铁性食物摄入）；铁剂治疗（在医生监督下合理补充铁剂）；防治结合的综合措施（发病率高、治疗容易但复发率高，需采取综合措施防控）。
\end{defbox}

% ----- 8.4 脊柱弯曲异常 -----
\section{脊柱弯曲异常（难点*）}

\begin{keybox}
\textbf{流行特点：}2019年后，国家卫生健康委疾控局正式将脊柱弯曲异常纳入全国学生常见病和健康影响因素监测与干预工作中。中国儿童少年脊柱弯曲异常初筛阳性率在2\%$\sim$3\%。

\textbf{发生原因：}
\begin{enumerate}[leftmargin=*]
    \item \imp{习惯性姿势}：不良站姿（身体重心习惯性偏向一边）、不良坐姿（坐立身体偏斜、歪头写字、过分靠近桌子）、不良走姿（上身左右晃动、双肩前倾、垂头含胸）
    \item \imp{桌椅高矮不适合}：桌椅高差太大或太小
    \item \imp{体力活动缺乏}："以静代动"的生活方式导致脊柱弯曲异常检出率显著上升
    \item \imp{营养和体质因素}：营养不足、体质孱弱的儿童少年，骨骼肌得不到充分发育，在同样条件下更容易发生脊柱弯曲异常
\end{enumerate}

\textbf{预防控制：}定期筛查并建立档案；建设健康环境和健康教育；创建良好政策环境、工作机制和防控平台。
\end{keybox}

% ----- 8.5 龋齿和牙周病 -----
\section{儿童少年龋齿和牙周病}

\begin{keybox}
\textbf{流行特点：}2010—2019年中国中小学生恒龋患率上升了6.4个百分点，女生恒龋患率始终高于男生。2015年第四次全国口腔流行病学调查结果显示，12岁儿童牙周健康率为41.6\%；牙龈出血检出率为58.4\%，牙石检出率为61.3\%。

\textbf{龋齿发生原因——四联致病因素模式：}
\begin{enumerate}[leftmargin=*]
    \item \imp{细菌和菌斑}
    \item \imp{宿主因素}：牙列不齐、釉质发育不良、抗酸能力弱等
    \item \imp{食物因素}
    \item \imp{时间因素}：龋齿是慢性硬组织破坏性疾病，菌斑在牙表面的滞留时间、菌斑内酸性产物的持续时间越长，发生龋齿的危险性越大
\end{enumerate}

\textbf{牙周病发生原因：}多疾病因素的综合作用结果，通常由细菌、宿主和环境三方面条件决定。

\textbf{预防控制：}
\begin{enumerate}[leftmargin=*]
    \item \imp{普及口腔保健，加强健康教育}：提高儿童和家长的口腔保健意识
    \item \imp{注意饮食卫生，增强宿主抗龋力}：限制糖的摄入对控制龋病有良好效果
    \item \imp{健全学校口腔疾病防治网}：大力推广氟化防龋、窝沟封闭等主动预防措施
\end{enumerate}
\end{keybox}

% ----- 复习题 -----
\section{复习题}

\begin{enumerate}[leftmargin=*, itemsep=8pt]
    \item \textbf{【名词解释】}近视；等效球镜（SE）；视力不良；高度近视
    \item \textbf{【名词解释】}单纯性肥胖；营养不足（生长迟缓、消瘦、体重不足）；四联致病因素
    \item \textbf{【简答题】}简述近视的流行特点、病因与危险因素。
    \item \textbf{【简答题】}简述户外活动预防近视的机制。
    \item \textbf{【简答题】}简述近视防控的综合策略（七项措施）。
    \item \textbf{【简答题】}简述儿童单纯性肥胖的流行趋势、病因及综合防控措施。
    \item \textbf{【简答题】}简述营养不足的三种表现形式及发生原因。
    \item \textbf{【简答题】}简述缺铁性贫血的发生原因及预防控制。
    \item \textbf{【简答题】}简述脊柱弯曲异常的发生原因及预防控制。
    \item \textbf{【简答题】}简述龋齿的四联致病因素模式及预防控制措施。
\end{enumerate}

\subsection*{参考答案要点}

\begin{enumerate}[leftmargin=*, itemsep=5pt]
    \item \textbf{近视}：人眼在调节放松状态下，平行光线经眼球屈光系统后聚焦在视网膜之前的屈光不正类型。\textbf{SE}：球镜度 + 1/2柱镜度，SE $\leq$ $-$0.50D为近视判定标准。\textbf{视力不良}：各种原因导致视力低于一定水平的总称，包括近视、远视、散光等。\textbf{高度近视}：SE $\leq$ $-$6.0D。

    \item \textbf{单纯性肥胖}：在遗传、环境因素交互作用下，因能量摄入超过能量消耗导致体内脂肪积聚过多的慢性代谢性疾病，与遗传、饮食、身体活动水平等有关。\textbf{营养不足}：包括生长迟缓（年龄别身高低）、消瘦（身高别体重低）、体重不足（年龄别体重低）。\textbf{四联致病因素}：龋齿由细菌（菌斑）、宿主（牙齿）、食物（糖）、时间四个因素共同作用。

    \item 流行特点：中国是世界上近视率最高的国家之一，检出率随年龄学龄上升，城市$>$乡村但差距缩小，女童$>$男童，低龄化趋势明显。病因：遗传因素（遗传度70\%$\sim$80\%）、环境因素（户外活动时间少是最重要危险因素）、体质健康因素（青春期突增阶段是近视程度加重关键期）。

    \item 户外较强光照促进视网膜多巴胺释放→抑制眼轴延长→延缓近视进展；开阔视野减轻眼调节负荷→放松睫状肌→预防和控制调节紧张性近视。\textbf{关键在于户外时间长短，而非运动类型}。

    \item 七项措施：(1)筛查视力不良与近视；(2)建立视力健康档案；(3)培养健康用眼行为；(4)建设视觉友好环境；(5)增加日间户外活动；(6)规范视力健康监测与评估；(7)分级诊疗与矫治（高危因素者改变行为→视力下降者医学验光→已近视者配镜→进展快者低浓度阿托品/OK镜）。

    \item 流行趋势：中国7$\sim$18岁肥胖检出率从1985年0.1\%增至2019年9.6\%（增长75.6倍），预计乡村将全面超过城市。病因：遗传（多基因遗传）、个体（能量过度摄取/消耗不足）、家庭（生命早期环境影响）、社会环境（"胖是健康"观念、膳食快餐化、媒体影响）。防控：强化家庭责任（科学饮食、积极活动、体重监测）、学校责任（营养课堂、改善供餐、保证活动时间）、医疗机构责任（孕期和儿童体重管理）、政府责任（知识普及、食品营销管理、完善体育设施）。

    \item 三种表现：生长迟缓（年龄别身高低，阻碍身体和认知潜力）、消瘦（身高别体重低，面临更高死亡风险）、体重不足（年龄别体重低）。发生原因：营养素摄入不足；疾病（肠道寄生虫、慢性疾病）；膳食结构和饮食习惯不良；早期生长潜力未充分发挥；体像认知和心理因素（盲目节食）。

    \item 原因：体内需求量大（生长突增、月经周期）；铁的额外丢失（月经不调性失血为首位原因、外伤、疾病、环境污染）。防控：去除病因；一般治疗（合理膳食、增加富铁食物）；铁剂治疗（医生监督下合理补充）；防治结合综合措施。

    \item 原因：习惯性姿势（不良站姿、坐姿、走姿）；桌椅高矮不适合；体力活动缺乏（"以静代动"）；营养和体质因素。防控：定期筛查并建立档案；建设健康环境和健康教育；创建良好政策环境、工作机制和防控平台。

    \item 四联因素：细菌（菌斑）—宿主（牙列不齐、釉质发育不良）—食物（糖）—时间（菌斑滞留和酸性产物持续时间越长危险性越大）。防控：普及口腔保健加强健康教育；注意饮食卫生增强宿主抗龋力；健全学校口腔疾病防治网（氟化防龋、窝沟封闭等）。
\end{enumerate}

\newpage
